% BT1157 signature-Clifford refinement insert.

% HOST-INDEPENDENT PREAMBLE (added by scripts/fix_insert_portability.py).
% This insert may be \input by either root manuscript, and they do not define the
% same environments or macros: w33_paper.tex has lemma/proposition/remark and
% \PSp/\Aut, photonic_holonet.tex has none of them. Define only what is missing,
% so the file is portable and re-running this is a no-op.
\makeatletter
\@ifundefined{theorem}{\newtheorem{theorem}{Theorem}}{}
\@ifundefined{lemma}{\newtheorem{lemma}[theorem]{Lemma}}{}
\@ifundefined{proposition}{\newtheorem{proposition}[theorem]{Proposition}}{}
\@ifundefined{remark}{\newtheorem{remark}[theorem]{Remark}}{}
\@ifundefined{corollary}{\newtheorem{corollary}[theorem]{Corollary}}{}
\@ifundefined{definition}{\newtheorem{definition}[theorem]{Definition}}{}
\makeatother
\providecommand{\PSp}{\mathrm{PSp}}
\providecommand{\Aut}{\mathrm{Aut}}
\providecommand{\spec}{\operatorname{spec}}


\begin{remark}[Scalar line versus orientation pseudoscalar]
\label{rem:signature-clifford-refinement}
The refined projective-sector calculation separates two roles that should not be
identified.  In the Boolean/Clifford completion,
\[
  16=1+15,\qquad 15=4+6+4+1.
\]
The extra $1$ is the scalar/vacuum empty-mask line.  The orientation pseudoscalar
is the grade-four line inside the nontrivial $15$-sector.  Therefore the K3
signature ledger $|\tau|=16=15+1$ should be read as projective $15$ plus scalar
completion, not as ``orientation bit outside the projective sector.''  The
orientation information is carried by the pseudoscalar line within the $15$.
\end{remark}
